\documentclass[spanish, 11pt]{exam}

%These tell TeX which packages to use.
\usepackage{array,epsfig}
\usepackage{amsmath}
\usepackage{amsfonts}
\usepackage{amssymb}
\usepackage{amsxtra}
\usepackage{amsthm}
\usepackage{mathrsfs}
\usepackage{color}
\usepackage{multicol}
\usepackage{verbatim}

\usepackage[utf8]{inputenc}
\usepackage[spanish]{babel}
\usepackage{eurosym}

\usepackage{graphicx}
\graphicspath{{../img/}}


%\printanswers
\nopointsinmargin
\pointformat{}

%Pagination stuff.
%\setlength{\topmargin}{-.3 in}
%\setlength{\oddsidemargin}{0in}
%\setlength{\evensidemargin}{0in}
%\setlength{\textheight}{9.in}
%\setlength{\textwidth}{6.5in}
%\pagestyle{empty}

\renewcommand{\solutiontitle}{\noindent\textbf{Sol:}\enspace}

\newcommand{\class}{4º Académicas}
\newcommand{\examdate}{\today}
\newcommand{\examnum}{Geometría Analítica}
\newcommand{\tipo}{A}


\newcommand{\timelimit}{50 minutos}



\pagestyle{head}
\firstpageheader{\includegraphics[width=0.2\columnwidth]{header_left}}{\textbf{Departamento de Matemáticas\linebreak \class}\linebreak \examnum}{\includegraphics[width=0.1\columnwidth]{header_right}}
\runningheader{\class}{\examnum}{Página \thepage\ of \numpages}
\runningheadrule

\begin{document}


\begin{questions}
\question Dada la recta $r\equiv 3x + y = 2$, halla una recta $s$, paralela a $r$, y otra perpendicular $t$, que pasen por el punto $P(2, – 1)$.
\begin{solution} \end{solution}

\question Halla el coeficiente $a$ para que la recta $ax + 4y = 11$ pase por el punto $P(1, 2)$
\begin{solution} \end{solution}

\question	Dados los siguientes vectores: $\overrightarrow u \left( {3,{\text{ }}2} \right){\text{ }}$ y ${\text{ }}\overrightarrow v \left( {1,{\text{ }}4} \right)$, calcula:
\begin{multicols}{4}
\begin{parts}
\part[] $\overrightarrow u \, + \,\overrightarrow v $ 
\begin{solution} $=\sqrt{156^2-65^2}=169$ \end{solution}
\part[]$\overrightarrow u \, - \,\overrightarrow v $  
\begin{solution}
%solve(156**2-169*x) 
144 y 25 \end{solution}

\part[]$2\overrightarrow u \, + 3\,\overrightarrow v $ 
\begin{solution} $h=64+225=289 \\ c_1^2=289\cdot64 \to c_1=136 \\
c_2^2=289\cdot 225 \to c_2=255$\end{solution}

\part[]$3\overrightarrow u \, - 4\,\overrightarrow v $ 
\begin{solution} $\to \sqrt{13^2-5^2}=12$ \end{solution}

\end{parts}
\end{multicols}

\question Averigua el punto simétrico de $A(5, -1)$ con respecto a $B(4, -2)$. 
\begin{solution} $c=\dfrac{20}{\tg45} = 20 \ mm $\end{solution}

\question Halla el punto medio del segmento de extremos $A(5, -1)$ y $B(4, -2)$ 
\begin{solution} $c=\dfrac{20}{\tg45} = 20 \ mm $\end{solution}

\question Dados los puntos $A(2, -3)$, $B(-1, 4)$ y $C(x, 3)$, determina el valor de x para que $A$, $B$ y $C$ estén alineados. 
\begin{solution} $c=\dfrac{20}{\tg45} = 20 \ mm $\end{solution}

\question Halla las coordenadas del vértice $D$ del paralelogramo $ABCD$, sabiendo que $A(-1, -2)$, $B(3, 1)$ y $C(1, 3)$.
\begin{solution} $c=\dfrac{20}{\tg45} = 20 \ mm $\end{solution}


\question Halla las ecuaciones paramétricas de la recta paralela a $2x - y + 3 = 0 $ y que pasa por el punto $P(4, 3)$.
\begin{solution} $c=\dfrac{20}{\tg45} = 20 \ mm $\end{solution}


\question Dadas las rectas:  $r \equiv \left\{ 
    {\begin{matrix}
	   {x = 2 - 4\lambda }  \\ 
	   {y =  - 2 + \lambda}  \\ 
    \end{matrix} } \right.$ 	y 	
    $s \equiv \left\{ {\begin{matrix}
   {x = 3 + 8\lambda }  \\ 
   {y =  - 1 - 2\lambda }  \\ 
\end{matrix} } \right.$  averigua su posición relativa. Si se cortan, di cuál es el punto de corte

\begin{solution} $c=\dfrac{20}{\tg45} = 20 \ mm $\end{solution}

\question ¿Cuál ha de ser el valor de $k$ para que estas dos rectas sean paralelas? $$x + 3y -2 = 0 \ \ \ \  kx + 2y + 3 = 0$$ 
\begin{solution} $c=\dfrac{20}{\tg45} = 20 \ mm $\end{solution}

\question Halla el valor de $k$ para que las rectas $2x - 3y + 4 = 0 $  $-3x + ky -1 = 0$ 
sean perpendiculares \begin{solution} $c=\dfrac{20}{\tg45} = 20 \ mm $\end{solution}



\question Dados los puntos A(-1, -1), B(1, 4) y C(5, 2), hallar 
\begin{multicols}{2}
\begin{parts}
\part[] Si están alineados   
%(sin(rad(60))-sin(rad(30)))/(sin(rad(60))+sin(rad(30)))
\begin{solution} $- \sqrt{3} + 2\approx0.267949192431123$\end{solution}
\part[] Mediana trazada desde B    
\begin{solution} $- 2 \sqrt{2} + 3\approx0.17157287525381$\end{solution}
\part[] Altura trazada desde A      
\begin{solution} $- \frac{\sqrt{3}}{4} + \frac{3}{4}\approx0.316987298107781$\end{solution}
\part[] Mediatriz del lado AB 
\begin{solution} $\frac{\sqrt{2} + \sqrt{6}}{2 \sqrt{2} + 2 \sqrt{3}}\approx0.614014407382354$\end{solution}
\end{parts}
\end{multicols}

\question Sean $A(1,0)$, $B(4,-3)$ y $C(5,2)$ los tres vértices de un triángulo. Hallar: 

\begin{parts}
\part[] La ecuación de la recta que pasando por A es paralela a la que pasa por B y C
\begin{solution} $- 2 \sqrt{2} + 3\approx0.17157287525381$\end{solution}
\part[] La ecuación de la mediana que pasa por C.     
\begin{solution} $- \frac{\sqrt{3}}{4} + \frac{3}{4}\approx0.316987298107781$\end{solution}
\end{parts}


\begin{comment}
\question 
\begin{multicols}{3}
\begin{parts}
\part[]  
\begin{solution} \end{solution}
\end{parts}
\end{multicols}
\end{comment}


\end{questions}
\end{document}

